\chapter{Introdución}
\hyphenation{In-te-re-se}

Obxectivos Xerais, Relación da Documentación que conforma a Memoria, Descrición do Sistema, Información Adicional de Interese (métodos, técnicas ou arquitecturas utilizadas, xustificación da súa elección, etc.).

El objetivo de este proyecto no es el de hacer la aplicación más 
eficiente de todas ni la más revolucionaria. 


A lo largo de toda esta etapa universitaria, se nos presentaron muchas formas 
de trabajo, muchas maneras de realizar una misma tarea. La eficiencia no debe ser la prioridad
de la vida, sino el disfrute de la misma, trabajando sobre esta mentalidad, la decisión de las tecnologías usadas fue clara:
se usaría Spring con Java para este proyecto. No hay tecnología que más amigable y disfrutable me pareciera en todas nuestras
enseñanzas.


El objetivo, entonces, es el demostrar que, aún sin utilizar las tecnologías más eficientes, es posible
crear algo completamente funcional. Siempre, claro, que midamos bien nuestro entorno de desarrollo. 
La eficiencia es importante, claro, pero no en todos los campos o proyectos.
Esta aplicación podría ser recreada de múltiples maneras haciendo uso de
cualquier otra tecnología o lenguaje, pero en nuestro caso fue utilizado Java con Spring
para el backend y JavaScript con React para el frontend. 